\documentclass[11pt,a4paper]{article}
\usepackage[T1]{fontenc}
\usepackage{lmodern,amsmath,amssymb,amsthm,mathtools}
\usepackage[margin=27mm]{geometry}
\usepackage[colorlinks=true,linkcolor=blue,urlcolor=blue,citecolor=blue]{hyperref}
\newtheorem{theorem}{Theorem}
\newtheorem{lemma}[theorem]{Lemma}
\newtheorem{corollary}[theorem]{Corollary}
\newtheorem{proposition}[theorem]{Proposition}
\newcommand{\Q}{\mathbb Q}\newcommand{\Z}{\mathbb Z}
\newcommand{\A}{\mathbb A}\newcommand{\Pj}{\mathbb P}
\DeclareMathOperator{\End}{End}
\DeclareMathOperator{\Tr}{Tr}
\DeclareMathOperator{\divisor}{div}
\title{Aggregate trace heights of elliptic multiples}
\author{}\date{8 November 2026}
\begin{document}\maketitle
\begin{abstract}
We prove a linear height comparison for the sum of the traces of
abscissas of nonzero multiples of the fixed elliptic point, and a
second comparison for the trace of their product when the labels
are odd. The main
geometric argument excludes all positive-dimensional abelian
translates through the specified points in a trace fiber, including
translates defined only over an algebraic extension. A theorem of
Von Buhren then supplies a linear comparison with an effective
geometric slope and an ineffective additive constant. Under the
arithmetic premises of the specified rational source, this proves
the existence of a defining polynomial with twenty-seven total
variables. The fixed coefficient needed to dominate the additive
constant is not determined. The result is therefore an existence
theorem, distinct from an explicit polynomial specification.
\end{abstract}

\section{Functions with one torsion zero and one pole}

Throughout the geometric argument the base field is algebraically
closed of characteristic zero. Let $E$ be an elliptic curve, let
$O$ be its identity, and write $\mathcal R=\End(E)$. This ring is
commutative and has no zero divisors. Let $T\in E$ have finite
order $N>1$. Suppose $f\in k(E)$ satisfies
\[
 \divisor(f)=m[T]-m[O]\qquad(m>0).
\]
Thus $f$ has its only zero at $T$ and its only pole at $O$.
Such a function exists when $m=N$, because $N(T-O)$ is a
principal divisor on an elliptic curve.

\begin{lemma}\label{lem:no-internal-cancellation}
Let $P\in E$ have infinite order. Let $n_1,\ldots,n_r$ be
nonzero integers prime to $N$, and let
$\alpha_1,\ldots,\alpha_r\in\mathcal R$. At every point $U$
of $E$, a pole of one function
$f([n_i]P+\alpha_iU)$ cannot coincide with a zero of another.
Consequently, if some $\alpha_i\ne0$, the product of these
functions has a pole and is nonconstant.
\end{lemma}
\begin{proof}
Suppose
\[
 [n_i]P+\alpha_iU=O,\qquad
 [n_j]P+\alpha_jU=T.
\]
The first equation requires $\alpha_i\ne0$; a constant factor
cannot be a zero or a pole because its argument is nontorsion.
Applying the two endomorphisms and subtracting gives
\[
 (n_i\alpha_j-n_j\alpha_i)P=-\alpha_iT.
\]
The right side is torsion. A nonzero endomorphism is an isogeny
and cannot take a nontorsion point to a torsion point. Hence
$n_i\alpha_j=n_j\alpha_i$.

Put $g=\gcd(|n_i|,|n_j|)$, $u=n_i/g$, and $v=n_j/g$.
There are integers $a,b$ with $au+bv=1$. Set
$\gamma=a\alpha_i+b\alpha_j$; the preceding equality gives
$\alpha_i=u\gamma$ and $\alpha_j=v\gamma$. Therefore
$R=[g]P+\gamma U$ satisfies $[u]R=O$ and $[v]R=T$.
The order of $R$ divides $|u|$, which is prime to $N$, whereas
$T$ has order $N>1$. This is impossible. Finally, every nonzero
$\alpha_i$ is surjective, so the corresponding factor has a
pole, and the absence of a zero at that point keeps it a pole
of the product.
\end{proof}

Two nontorsion points $P_1,P_2$ will be called independent over
$\mathcal R$ modulo torsion if
$\alpha P_1+\beta P_2$ torsion implies $\alpha=\beta=0$.

\begin{lemma}\label{lem:cross-cancellation}
Suppose $P_1,P_2$ are independent over $\mathcal R$ modulo
torsion. If $n,m$ are nonzero integers and $\alpha,\beta$ are
nonzero endomorphisms, the functions
$f([n]P_1+\alpha U)$ and $f([m]P_2+\beta U)$ have no common
pole.
\end{lemma}
\begin{proof}
A common pole would give $\alpha U=-[n]P_1$ and
$\beta U=-[m]P_2$. Commutativity then gives
$n\beta P_1-m\alpha P_2=O$, contradicting independence.
\end{proof}

\section{Excluding abelian translates in an aggregate fiber}

Fix nonzero scalars $c_1,c_2$ and an integer $r\ge1$. On
$E^{2r}$ consider the rational function
\[
 F(Q_1,\ldots,Q_r,Q'_1,\ldots,Q'_r)
 =c_1\prod_{i=1}^r f(Q_i)+c_2\prod_{i=1}^r f(Q'_i).
\]

\begin{theorem}\label{thm:no-translate}
Let $P_1,P_2$ be independent over $\mathcal R$ modulo torsion.
Suppose every $n_i,m_i$ is nonzero and prime to $N$. Put
\[
 q=([n_1]P_1,\ldots,[n_r]P_1,
       [m_1]P_2,\ldots,[m_r]P_2).
\]
No positive-dimensional abelian translate through $q$ is
contained in a level set of $F$.
\end{theorem}
\begin{proof}
Suppose $q+B$ is such a translate, with $B$ a nonzero abelian
subvariety of $E^{2r}$. There is a nonzero homomorphism
$\phi:E\to B$. For example, a nonzero coordinate projection
$B\to E$ is surjective; dualizing and composing with an inverse
to a polarization up to multiplication produces such a map.
Write its coordinates as
$(\alpha_1,\ldots,\alpha_r,\beta_1,\ldots,\beta_r)$,
not all zero. The restriction of $F$ to $q+\phi(E)$ is
\[
 c_1\prod_i f([n_i]P_1+\alpha_iU)
 +c_2\prod_i f([m_i]P_2+\beta_iU).
\]
Within either product a pole survives by
Lemma~\ref{lem:no-internal-cancellation}. A pole of one product
is not a pole of the other by
Lemma~\ref{lem:cross-cancellation}. Since at least one
endomorphism is nonzero, at least one product has a pole, and
it remains a pole of the sum. This contradicts constancy.
\end{proof}

The same proof permits any finite number of blocks based on
points that are pairwise independent over $\mathcal R$ modulo
torsion. The integers in each block need only be prime to the
order of the unique torsion zero of its chosen function.

\section{Sums of abscissas without a parity restriction}

Give $E$ a Weierstrass equation and write $x$ for its abscissa.
This function has its only pole at $O$, of order two. Let
$w_1,\ldots,w_r,w'_1,\ldots,w'_r$ be positive rational numbers
and let $c_1,c_2$ be nonzero algebraic numbers. Define
\[
 F_+(Q_1,\ldots,Q_r,Q'_1,\ldots,Q'_r)
   =c_1\sum_i w_i x(Q_i)+c_2\sum_i w'_i x(Q'_i).
\]

\begin{lemma}\label{lem:sum-poles}
Let $P\in E$ have infinite order, let all integers $n_i$ be
nonzero, and let $\alpha_i\in\mathcal R$. If some
$\alpha_i\ne0$, the function
\[
 \sum_i w_i x([n_i]P+\alpha_iU)
\]
has a pole and is nonconstant.
\end{lemma}
\begin{proof}
Choose a pole $U_0$ of a nonconstant summand, say the $i$th.
If another summand, the $j$th, has a pole there, then
$[n_i]P+\alpha_iU_0=[n_j]P+\alpha_jU_0=O$.
Applying the endomorphisms and subtracting gives
$(n_i\alpha_j-n_j\alpha_i)P=O$ and hence
$n_i\alpha_j=n_j\alpha_i$.

Let $z$ be a uniformizer at $O$ with
$x=z^{-2}+O(z^{-1})$, and use $z(U-U_0)$ locally at $U_0$.
Let $\lambda_j$ be the nonzero tangent multiplier of
$\alpha_j$ at $O$. The leading coefficient of the $j$th
summand at this pole is $w_j\lambda_j^{-2}$, and the
endomorphism identity gives
$\lambda_j=(n_j/n_i)\lambda_i$. Consequently the coefficient
of $z^{-2}$ in the full sum is
\[
 \lambda_i^{-2}
       \sum_{j:\,[n_j]P+\alpha_jU_0=O}
                    w_j(n_i/n_j)^2.
\]
The sum is a nonzero positive rational number. The pole
therefore cannot cancel.
\end{proof}

\begin{theorem}\label{thm:sum-no-translate}
Suppose $P_1,P_2$ are independent over $\mathcal R$ modulo
torsion, and all $n_i,m_i$ are nonzero. No
positive-dimensional abelian translate through
\[
 q=([n_i]P_1)_{i=1}^r\times([m_i]P_2)_{i=1}^r
\]
is contained in a level set of $F_+$.
\end{theorem}
\begin{proof}
Restrict to an elliptic subgroup direction exactly as in
Theorem~\ref{thm:no-translate}. Every nonconstant block sum
has a pole by Lemma~\ref{lem:sum-poles}. A pole in one block
cannot occur in the other block, because the same argument
as Lemma~\ref{lem:cross-cancellation} would give a nontrivial
endomorphism relation between $P_1$ and $P_2$. At least one
block is nonconstant, so their weighted sum is nonconstant.
\end{proof}

This result uses no torsion zero of $x$ and imposes no parity
condition. Its positivity assumption concerns the fixed
rational weights within each block, not the values of the
abscissas.

\begin{theorem}[Several blocks]\label{thm:many-blocks}
Let $P_1,\ldots,P_s$ be nontorsion points on $E$ such that
each pair is independent over $\mathcal R$ modulo torsion.
Fix positive integers $r_j$, positive rational weights
$w_{ji}$, and nonzero algebraic scalars $c_j$. On
$E^{\sum_jr_j}$ put
\[
 F_{\mathrm{sum}}((Q_{ji}))
      =\sum_{j=1}^s c_j\sum_{i=1}^{r_j}w_{ji}x(Q_{ji}).
\]
For any tuple $q=([n_{ji}]P_j)$ with all $n_{ji}\ne0$, no
positive-dimensional abelian translate through $q$ is
contained in a fiber of $F_{\mathrm{sum}}$.

There is a product version. For each $j$ suppose
$\divisor(f_j)=m_j[T_j]-m_j[O]$, where $T_j$ has order
$N_j>1$. Replace $F_{\mathrm{sum}}$ by
\[
 F_{\mathrm{product}}((Q_{ji}))
      =\sum_{j=1}^s c_j\prod_{i=1}^{r_j}f_j(Q_{ji}).
\]
The same conclusion holds if every $n_{ji}$ is prime to $N_j$.
\end{theorem}
\begin{proof}
Restrict to a nonzero homomorphism $E\to B$ in any proposed
positive-dimensional subgroup $B$. Some block has a
nonconstant coordinate. Within that block a pole survives
by Lemma~\ref{lem:sum-poles}, or by
Lemma~\ref{lem:no-internal-cancellation} in the product case.
A pole cannot occur in two distinct blocks: their two pole
equations would contradict the pairwise independence of
their base points. The pole therefore survives in the full
sum, contradicting constancy.
\end{proof}

The scalars $c_j$ need not be real or positive. Pairwise
independence suffices; no independence assumption on the
entire collection is required.

\section{A linear height comparison}

We use the following consequence of Von Buhren's family height
theorem \cite[Theorem~1.1]{VonBuhren}. Let $A$ be an abelian
variety, $S$ a projective variety, and $V\subseteq A\times S$
an integral closed subvariety. For a fixed subgroup
$\Gamma\subseteq A(\overline\Q)$ of finite rank, there are
constants $\mu,\nu\ge0$ such that
\begin{equation}\label{eq:VB}
 \widehat h_A(a)\le\mu h_S(s)+\nu
\end{equation}
whenever $a\in\Gamma$ lies in the fiber $V_s$ and does not lie
on a positive-dimensional abelian translate contained in that
fiber. The slope $\mu$ can be computed from projective
geometric data. The theorem does not provide an effective
value for $\nu$.

\begin{corollary}\label{cor:height}
Assume the data of Theorem~\ref{thm:no-translate} are algebraic
numbers. There are constants $M>0$ and $M_0\ge0$, independent
of all the integers $n_i,m_i$, such that
\[
 \sum_{i=1}^r\bigl(\widehat h_E([n_i]P_1)
                 +\widehat h_E([m_i]P_2)\bigr)
 \le M h(F(q))+M_0
\]
for every tuple in that theorem. Here $h$ is the absolute
logarithmic height on $\Pj^1$. One may choose an effective
geometric slope $M$ after fixing compatible height
normalizations. No effective $M_0$ is asserted.
The same conclusion holds with $F_+$ and the hypotheses of
Theorem~\ref{thm:sum-no-translate}, allowing all nonzero labels.
\end{corollary}
\begin{proof}
Let $A=E^{2r}$ and take for $V\subseteq A\times\Pj^1$ the
closure of the graph of $F$. This is integral and projective.
Let $\Gamma$ be the subgroup generated by $P_1$ in each of
the first $r$ coordinates and by $P_2$ in each of the last
$r$ coordinates. It is finitely generated.

Every indicated $q$ is in the domain of $F$, since no
$[n_i]P_1$ or $[m_i]P_2$ is the identity. If an abelian
translate through $q$ were contained in the projective fiber
$V_{F(q)}$, it would meet the domain of $F$ at $q$. Its dense
open intersection with that domain would make $F$ constant
on the translate. Theorem~\ref{thm:no-translate} rules this
out. This argument also handles possible boundary components
introduced by closing the graph; none through $q$ changes
the conclusion.

Apply~\eqref{eq:VB}, using a symmetric ample product line
bundle on $A$. Its canonical height is a fixed positive
multiple of the sum of the coordinate canonical heights.
Rescaling gives the stated inequality. For $F_+$ the same
proof uses Theorem~\ref{thm:sum-no-translate} to exclude the
exceptional translates.
\end{proof}

For either function of Theorem~\ref{thm:many-blocks} the same
argument gives a uniform inequality
\[
 \sum_{j,i}\widehat h_E([n_{ji}]P_j)
       \le M h(F(q))+M_0.
\]
The subgroup is generated by the finitely many base points
placed in their respective coordinates. The constants are
fixed before any of the integer labels are chosen.

\section{The quadratic-field elliptic curve}

Let
\[
 K=\Q(\sqrt2),\quad d=3-\sqrt2,\quad d'=3+\sqrt2,
 \quad E:\ y^2=x^3-d^2x,
\]
and let $\sigma$ be the nonidentity automorphism of $K/\Q$.
The point used in the rational source is
\[
 R_0=(1+\sqrt2,2\sqrt2),\qquad
 P=2R_0=\left(\frac{57}{8}-\frac72\sqrt2,
             -\frac{165}{8}+\frac{417}{32}\sqrt2\right).
\]
For this section its infinite order is the only group
property needed; no assertion about the full Mordell--Weil
group is required.

\begin{lemma}
The displayed point $P$ has infinite order.
\end{lemma}
\begin{proof}
Substitution verifies $R_0\in E(K)$, and the usual tangent
formula gives the displayed double. At the prime
$\mathfrak p=(\sqrt2)$ normalize $v_{\mathfrak p}(\sqrt2)=1$,
so $v_{\mathfrak p}(2)=2$. The displayed coordinates give
\[
 v_{\mathfrak p}(x(P))=-6,\qquad
 v_{\mathfrak p}(y(P))=-9,\qquad
 v_{\mathfrak p}(-x(P)/y(P))=3.
\]
The parameter $z=-x/y$ places $P$ in the formal group of
this integral Weierstrass equation. Its multiplication series
have integral coefficients and linear term $[m](z)=mz+O(z^2)$.
For odd $m$ and $v(z)\ge3$, the linear term uniquely has
least valuation, so $v([m]z)=v(z)$. For multiplication by
two, the linear term has valuation $v(z)+2<2v(z)$,
and therefore $v([2]z)=v(z)+2$. No successive multiplication
by an odd integer and by a power of two can give zero.
Every nonzero integer has this form, proving the claim.
\end{proof}

\begin{lemma}\label{lem:independent}
Over a field containing $\sqrt d$ and $\sqrt{d'}$, identify
$E$ and $E^\sigma$ with $E_0:Y^2=X^3-X$ by
\[
 (X,Y)=\left(\frac{x}{d},\frac{y}{d\sqrt d}\right),
 \qquad
 (X',Y')=\left(\frac{x'}{d'},\frac{y'}{d'\sqrt{d'}}\right).
\]
The images $P_1,P_2$ of $P,P^\sigma$ are independent over
$\End(E_0)=\Z[i]$ modulo torsion.
\end{lemma}
\begin{proof}
The squareclasses of $d,d'$ are independent in $K(i)$.
Indeed, a positive element of the real field $K$ that is a
square in $K(i)$ is already a square in $K$: writing a square
root as $a+ib$ gives $ab=0$, and the purely imaginary case
would be negative at a real embedding. Neither $d$ nor $d'$
is a square in $K$, since its norm to $\Q$ is $7$. Their
product $7$ is not a square in $K$: expanding
$(a+b\sqrt2)^2=7$ for rational $a,b$ gives $ab=0$ and then
requires either $7$ or $7/2$ to be a rational square.

It follows that $K(i,\sqrt d,\sqrt{d'})/K(i)$ has an
automorphism fixing $\sqrt d$ and negating $\sqrt{d'}$.
It fixes $P_1$, negates $P_2$, and fixes every endomorphism
in $\Z[i]$. If $\alpha P_1+\beta P_2$ were torsion, applying
the automorphism and subtracting would make $2\beta P_2$
torsion. Infinite order forces $\beta=0$, and then
$\alpha=0$.
\end{proof}

\begin{theorem}\label{thm:aggregate}
For any fixed $r\ge1$, there are constants $M_+>0$ and
$M_{+,0}\ge0$ such that, for all nonzero integers
$n_1,\ldots,n_r$,
\begin{equation}\label{eq:sum-height}
 \sum_{i=1}^r\widehat h_E([n_i]P)
 \le M_+ h\!\left(\frac12\Tr_{K/\Q}
                 \sum_{i=1}^r x([n_i]P)\right)+M_{+,0}.
\end{equation}
There are also constants $M_\times>0$, $M_{\times,0}\ge0$
such that, when all the labels are odd,
\begin{equation}\label{eq:aggregate-height}
 \sum_{i=1}^r\widehat h_E([n_i]P)
 \le M_\times h\!\left(\frac12\Tr_{K/\Q}
                 \prod_{i=1}^r x([n_i]P)\right)+M_{\times,0}.
\end{equation}
Both slopes can be chosen effectively from the fixed geometry;
the additive constants are not determined effectively here.
\end{theorem}
\begin{proof}
On $E_0$, the coordinate $X$ has divisor
$2[(0,0)]-2[O]$. Its zero is a nonzero point of order two,
so the index condition in Theorem~\ref{thm:no-translate} is
exactly oddness. Under the isomorphisms of
Lemma~\ref{lem:independent}, the rational function in the
height on the right of~\eqref{eq:aggregate-height} is
\[
 F=\frac{d^r}{2}\prod_{i=1}^r X_i
       +\frac{(d')^r}{2}\prod_{i=1}^r X'_i.
\]
For the sum, the corresponding function is
$F_+=(d/2)\sum_iX_i+(d'/2)\sum_iX'_i$ and
Theorem~\ref{thm:sum-no-translate} applies without parity
restrictions. Apply Corollary~\ref{cor:height} to both
functions with $m_i=n_i$.
Canonical heights are unchanged by these pointed
isomorphisms, and conjugation preserves absolute canonical
height. The coordinate height sum there is twice the sum
in~\eqref{eq:aggregate-height}; divide by two.
\end{proof}

Notice that the inequality is linear. Finiteness of each
individual fiber alone would not give this conclusion. The
additive constant is the remaining effectiveness issue; it is
not an additional mathematical conjecture.

\section{A trace-height consequence over number fields}

\begin{corollary}\label{cor:number-fields}
Let $K/F$ be a finite extension of number fields, let $E/K$
be an elliptic curve, and let $P\in E(K)$ have infinite order.
For every embedding $\sigma:K\hookrightarrow\overline F$
over $F$, suppose there is a pointed isomorphism
$\phi_\sigma:E^\sigma\to E_0$ to a fixed elliptic curve over
$\overline F$. Assume that the points
$P_\sigma=\phi_\sigma(P^\sigma)$ are pairwise independent
over $\End(E_0)$ modulo torsion.

For every fixed positive integer $r$ there are constants
$M>0$, $M_0\ge0$ such that, for all nonzero integer labels,
\[
 \sum_{i=1}^r\widehat h_E([n_i]P)
 \le M h\!\left(\Tr_{K/F}
                    \sum_{i=1}^r x([n_i]P)\right)+M_0.
\]
The slope can be chosen effectively from the fixed geometric
data; the additive constant has the same ineffective status
as in Corollary~\ref{cor:height}.
\end{corollary}
\begin{proof}
Under a pointed isomorphism of Weierstrass models the
abscissa transforms as
$x^\sigma=a_\sigma x_0+b_\sigma$, with $a_\sigma\ne0$.
The function in the height is therefore
\[
 \sum_\sigma a_\sigma\sum_i x_0([n_i]P_\sigma)
                  +r\sum_\sigma b_\sigma.
\]
Theorem~\ref{thm:many-blocks} excludes constant translates
through the indicated points for the nonconstant part.
Adding the fixed scalar does not change that conclusion.
The graph-family argument gives a linear bound for the sum
of the canonical heights in all the embedding blocks.
Absolute canonical height is unchanged by conjugation and
pointed isomorphism, so this sum is $[K:F]$ times the left
side of the displayed inequality. Divide by $[K:F]$.
\end{proof}

This corollary is a height comparison under explicit
independence hypotheses. It does not assert a Diophantine
definition of the ring of integers of $K$.

\section{A conditional existence result with twenty-seven variables}

We specify precisely how the height theorem applies to the
explicit rational source in~\cite{Baseline29}. We retain its group, contact-height,
pole, and simultaneous prime-pattern premises. Its contact
exponent is $\eta=19$, its comparison orders are $39$ and
$10$, and its scalar order is $20$. These premises and the
same finite-support completeness argument are the only
arithmetic inputs in the following consequence.

Write the four root labels as $(n,a,t,b)$, all nonzero integers, and put
\[
 B=\max(|n|,|a|,|t|,|b|),\quad c=|a-1|,
 \quad\Delta=t^2-a(a-2)tb+(a-2)^2b^2-1,
\]
and use the positive rational number
\[
 F=\Tr_{K/\Q}
       \bigl(x([n]P)+x([a]P)+x([t]P)+x([b]P)\bigr).
\]
Positivity follows because the source points are nonzero
 doubles and their abscissas are positive at both real
embeddings. In particular $F\ne0$. The sum uses only additions
after the four abscissas have been computed. Replacing the
half trace in Theorem~\ref{thm:aggregate} by the full trace
changes the height by at most $\log 2$, absorbed into its
additive constant. The product theorem gives an independent
height comparison when all labels are odd; the construction
here uses the sum.

Use only the two comparison scales $1$ and $F^2$. The
source's contact-height inequality for these scales gives
\begin{equation}\label{eq:two-scales}
 h(F)\le(H+\log3)|N|^{19/k}
\end{equation}
whenever their local certificates give the required
$k$-fold contact divisibility for the nonzero integer $N$.
Indeed, the unscaled and scaled arguments bound the heights
of $\lambda$ and $F^2\lambda$ by the same right side, so
$2h(F)=h(F^2)\le h(\lambda)+h(F^2\lambda)$.
Theorem~\ref{thm:aggregate}, with $r=4$, then implies
\begin{equation}\label{eq:B-height}
 B^2\le M_1 h(F)+M_2
\end{equation}
for fixed $M_1>0$, $M_2\ge1$. This follows directly from
$\widehat h_E([j]P)=j^2\widehat h_E(P)$ and
$\widehat h_E(P)>0$.

\begin{theorem}\label{thm:27}
Under the stated arithmetic premises, there exists a polynomial
over $\Z$ in twenty-seven total variables, including its input
variable, of total degree at most $2^{158}$, whose rational
zero set projects exactly to $\Z$.
The proof does not determine all of its coefficients.
\end{theorem}
\begin{proof}
For any positive odd integer $L$, the scalar normalizer can be
chosen as
\[
 R_L(X)=\frac{102^2\ell^L r(X)}{(1+C_Lr(X))^2},
 \qquad \ell=76249,\qquad
 C_L=17\bigl(\ell^{(L+1)/2}+1\bigr),
\]
with $\beta=216X^4+450459$ and
$r=\beta(\beta-1)(\beta-2401)(32\beta+343)(2\beta-1)$.
The finite-field certificate in~\cite{ScalarNormalization}
proves that $r$ and $1+17r$ have no
zero modulo $\ell$. Thus
\[
 v_\ell(R_L)=L\quad(v_\ell X\ge0),\qquad
 v_\ell(R_L)=L-20v_\ell X\quad(v_\ell X<0).
\]
Also $0<R_L<1$, since
$C_L^2r(0)>17^2\ell^{L+1}r(0)>102^2\ell^L$.
All other scalar local conditions and square-denominator
conditions are unchanged. The fixed-prime extraction argument
therefore gives
\[
 v_\ell(a-1)\ge19L+1,
 \qquad c\ge\ell^{19L+1}.
\]

If $\Delta\ne0$, local comparison and~\eqref{eq:two-scales}
with $k=39$ give, using $|\Delta|<5B^4$,
\[
 B^2\le M_1(H+\log3)(5B^4)^{19/39}+M_2.
\]
The exponent $76/39$ is strictly smaller than two. Hence this
inequality fails for all sufficiently large $B$. Since
$B\ge|a|\ge c-1$, choosing the fixed odd $L$ sufficiently
large excludes $\Delta\ne0$ for every solution.

Now $\Delta=0$, and the elementary integer Lucas growth
argument gives $\log B>c-4$ once $|a|>7$.
If $\beta\ne n$, write $N$ for the absolute numerator of
$\beta-n$ in lowest terms. The source gives
$N\le2e^{h(\beta)}B$ and contact divisibility of order ten.
Equations~\eqref{eq:two-scales} and~\eqref{eq:B-height} give
\[
 B^2\le M_1(H+\log3)
               (2e^{h(\beta)}B)^{19/10}+M_2.
\]
For $B^2>2M_2$, taking logarithms yields
$\log B\le19h(\beta)+M_3$ for a fixed constant $M_3$.
The scalar estimate remains $h(\beta)\le Hc^{19/20}$.
It follows that
\[
 c-4<19Hc^{19/20}+M_3,
\]
which fails for all sufficiently large $c$. Increase the
same fixed odd $L$ if necessary. We conclude $\beta=n\in\Z$,
and the source input map then gives $X\in\Z$.

Conversely, for each integer $X$ the unchanged completeness
argument prepares the finite support of $R_L(X)$ and applies
the simultaneous prime-pattern premise to the two remaining
scales. Removing a comparison scale removes requirements;
the source's finite-avoidance argument still applies. It
therefore supplies all rational witnesses.

The retained coordinates are the input, eighteen coordinates
for nine elliptic point charts, one slope, four parameters for
the two comparison poles, and two parameters for the scalar
pole: $1+18+1+4+2=26$. The finite-algebra elimination
theorem of~\cite{FiniteAlgebra} adds one coordinate, giving twenty-seven. Its rational
projection is unchanged by that elimination. Clearing the
fixed rational coefficients gives a polynomial over $\Z$.

All constants used to choose $L$ exist by the proved height
theorem. Their existence proves the polynomial's existence.
The ineffective additive constant $M_2$ prevents this proof
from specifying a certified numerical value of $L$.

The direct point charts and the single cubic change of basis
for their square roots are proved in~\cite{SquareCharts,QuadraticCharts}.
The resulting finite-algebra specification
\cite{FamilyCertificate} has nine quartic point generators
and three quadratic pole generators. Its maximum rank is
$4^9\cdot2^3=2^{21}$, with eight source cases. Exact degree
arithmetic gives the majorant
\[
 232075682770882905150152990100588354340946706432
                 <2^{158}.
\]
The exponent $L$ enters only integer coefficients. Every
degree computation and the same majorant are therefore
uniform in the fixed odd $L\ge3$, including every
sufficiently large member used above. This proves the degree
assertion without an effective value of $L$.
\end{proof}

The formal logarithm orders remain $39$, $10$, and $20$;
only a fixed scalar coefficient is enlarged. Thus the
ineffectiveness does not conceal an extra coordinate or a
degree that grows with $L$: $L$ occurs solely in coefficients.
The explicit twenty-nine-variable construction
remains a different deliverable, with all constants given.

\begingroup\small
\begin{thebibliography}{9}
\setlength{\itemsep}{2pt}
\setlength{\parskip}{0pt}
\bibitem{VonBuhren}
J.~Von Buhren, \emph{Borne de hauteur semi-effective pour le
probl\`eme de Mordell--Lang dans une vari\'et\'e ab\'elienne},
Journal de Th\'eorie des Nombres de Bordeaux
\textbf{29} (2017), no.~1, 307--320.
\href{https://doi.org/10.5802/jtnb.981}{doi:10.5802/jtnb.981}.
\href{https://www.numdam.org/item/10.5802/jtnb.981.pdf}{Full text}.

\bibitem{Baseline29}
\emph{An explicit polynomial in 29 variables over the rationals},
project manuscript, 8 November 2026.
\href{https://github.com/JoseBrox/diophantine-z-q_improvements}{Project repository};
source file \path{few_variables_20261108_v8.tex}.

\bibitem{FiniteAlgebra}
\emph{Eliminating finite algebraic choices with one coordinate},
project manuscript, 7 November 2026.
\href{https://github.com/JoseBrox/diophantine-z-q_improvements}{Project repository};
source file \path{finite_algebra_20261107_v7.tex}.

\bibitem{ScalarNormalization}
\emph{Scalar contact normalization by finite-field avoidance},
companion manuscript, 8 November 2026.
Source file \path{scalar_normalization_20261108_v9.tex};
exact certificates and independent checker accompany the manuscript.

\bibitem{FamilyCertificate}
\emph{A defining polynomial in twenty-seven variables:
an existence theorem with a uniform degree bound},
polynomial specification with independent verification,
9 November 2026; source file
\path{few_variables27_nontorsion_20261109_v12.tex}.
The reproducibility archive contains the polynomial
specification, construction report, source verification and
independent verification for the displayed example coefficient
$L=137731$. The example is used to verify the uniform degree
calculation; it is not asserted to exceed the ineffective
sufficient threshold.
\bibitem{SquareCharts}
\emph{Square abscissas and direct charts for the elliptic subgroup},
companion manuscript, 8 November 2026.
Source file \path{square_abscissa_charts_20261108_v9.tex}.

\bibitem{QuadraticCharts}
\emph{A single quartic chart for square roots in a quadratic extension},
companion manuscript, 8 November 2026.
Source file \path{quadratic_square_chart_20261108_v9.tex}.
\end{thebibliography}
\endgroup
\end{document}
